%<<>>==<<>><<>>==<<>><<>>==<<>><<>>==<<>><<>>==<<>>
\documentclass[aspectratio=169,10pt]{beamer}
%<<>>==<<>><<>>==<<>><<>>==<<>><<>>==<<>><<>>==<<>>
\usepackage[utf8]{inputenc}
\usepackage[T1]{fontenc}
\usepackage[spanish,mexico]{babel}
\usepackage{amsmath,amssymb,mathtools}
\usepackage{booktabs,array,multicol}
\usepackage{xcolor,graphicx,tikz,pgfplots}
\pgfplotsset{compat=1.18}
\pgfplotsset{unbounded coords=jump}
%<<>>==<<>><<>>==<<>><<>>==<<>><<>>==<<>><<>>==<<>>
\usetheme{Madrid}
\usecolortheme{default}
\definecolor{azul}{RGB}{25,70,125}
\setbeamercolor{structure}{fg=azul}
\setbeamercolor{frametitle}{bg=azul,fg=white}
\setbeamertemplate{navigation symbols}{}
\setbeamertemplate{footline}{%
  \leavevmode\hbox{%
  \begin{beamercolorbox}[wd=.28\paperwidth,ht=2.5ex,dp=1ex,leftskip=1ex]{author in head/foot}Carlos Ernesto Martínez Rodríguez\end{beamercolorbox}%
  \begin{beamercolorbox}[wd=.52\paperwidth,ht=2.5ex,dp=1ex,center]{title in head/foot}Cálculo Diferencial -- Unidad 2: Funciones\end{beamercolorbox}%
  \begin{beamercolorbox}[wd=.20\paperwidth,ht=2.5ex,dp=1ex,rightskip=1ex plus1fil]{date in head/foot}\insertframenumber/\inserttotalframenumber\end{beamercolorbox}}}
  %<<>>==<<>><<>>==<<>><<>>==<<>><<>>==<<>><<>>==<<>>
\newtheorem{definicion}{Definición}
\newtheorem{teorema}{Teorema}
\newtheorem{proposicion}{Proposición}
\newtheorem{corolario}{Corolario}
\newtheorem{observacion}{Observación}
\newtheorem{ejemplo}{Ejemplo}
\newtheorem{nota}{Nota}
%<<>>==<<>><<>>==<<>><<>>==<<>><<>>==<<>><<>>==<<>>
\newcommand{\R}{\mathbb{R}}
\newcommand{\Dom}{\operatorname{Dom}}
\newcommand{\Ima}{\operatorname{Im}}
\newcommand{\Ran}{\operatorname{Ran}}
\newcommand{\id}{\operatorname{id}}
%<<>>==<<>><<>>==<<>><<>>==<<>><<>>==<<>><<>>==<<>>
\title{Cálculo Diferencial\\Unidad 2: Funciones}
\subtitle{Dominio, Imagen, Operación y composición de Funciones}
\author{Carlos Ernesto Martínez Rodríguez}
\institute{Tecnológico Nacional de México}
\date{}
%<<>>==<<>><<>>==<<>><<>>==<<>><<>>==<<>><<>>==<<>>
\begin{document}
%<<>>==<<>><<>>==<<>><<>>==<<>><<>>==<<>><<>>==<<>>

\begin{frame}{Composición de funciones}
Determinar $f\circ g$, $g\circ f$, $f\circ f$ y $g\circ g$

Para $f(x)=3-2x$; $g(x)=6-3x$, calcular
\begin{itemize}
\item $f(3a)$: $f(3a)= 3-2a$
\item $f(\frac{3}{a})$: $f(\frac{3}{a})= 3-2(\frac{3}{a})$; 
\item $f(\frac{1}{2})$: $f(\frac{1}{2})= 3-2(\frac{1}{2})$
\item $f(\frac{1}{2}x)$: $f(\frac{1}{2}x)= 3-2(\frac{1}{2}x)=3-x$
\item $f(-3x)$: $f(-3x)= 3-2(-3x)=3+6x$
\item $f(6-3x)$: $f(6-3x)= 3-2(6-3x)=3-12+6x=-9+6x$
\item  $f(g(x))$: $f(g(x))= f(6-3x)=3-2(6-3x)= -9+6x$
\item  $f(f(x))$: $f(f(x))= f(3-2x)=3-2(3-2x)=3-6+4x=-3+4x$
\end{itemize}

$Dominio(f)=\mathbb{R}$

$Dominio(f\circ g)=\mathbb{R}$

$Dominio(f\circ f)=\mathbb{R}$
\end{frame}

\begin{frame}{Composición de funciones}
Determinar $f\circ g$, $g\circ f$, $f\circ f$ y $g\circ g$

Para $f(x)=3-2x$; $g(x)=6-3x$, calcular
\begin{itemize}
\item $g(3a)$: $g(3a)= 6-3(a)=6-3a$
\item $g(\frac{3}{a})$: $g(\frac{3}{a})= 6-3(\frac{3}{a})=6-\frac{9}{a}$
\item $g(\frac{1}{2})$: $g(\frac{1}{2})= 6-\frac{3}{2}=\frac{9}{2}$
\item $g(\frac{1}{2}x)$: $g(\frac{1}{2}x)=6-3(\frac{1}{2}x)=6-\frac{3}{2}x$
\item $g(-2x)$: $g(-2x)= 6-3(-2x)=6+6x$
\item $g(3-2x)$: $g(3-2x)= 6-3(3-2x)=6-9+6x=-3+6x$
\item  $g(f(x))$: $g(f(x))= g(3-2x)=6-3(3-2x)=6-9+6x=-3+6x$
\item $g(-3x)=6-3(-3x)=6+9x$
\item $g(6-3x)=6-3(6-3x)=6-18+9x=-12+9x=g(g(x))$
\end{itemize}
$Dominio(g)=\mathbb{R}$

$Dominio(g\circ f)=\mathbb{R}$

$Dominio(f\circ g)=\mathbb{R}$



\end{frame}

\begin{frame}{Composición de funciones}
Determinar $f\circ g$, $g\circ f$, $f\circ f$ y $g\circ g$

Para $f(x)=3-2x$; $g(x)=6-3x$, calcular

$Dominio(f)=\mathbb{R}$

$Dominio(f\circ g)=\mathbb{R}$

$Dominio(f\circ f)=\mathbb{R}$

$Dominio(g)=\mathbb{R}$

$Dominio(g\circ f)=\mathbb{R}$

$Dominio(f\circ g)=\mathbb{R}$

\end{frame}

\begin{frame}
Defina las siguientes funciones y determine el dominio de las funciones
resultantes: (a) $f(x^2)$; (b) $[f(x)]^2$; (c) $(f\circ f )(x)$; (d) $(f\circ f)(-x)$, para  $f(x) = \frac{1}{x-1}$

\begin{itemize}

\item $f(x^{2})=\frac{1}{x^2-1}=\frac{1}{(x-1)(x+1)}$

\item $[f(x)]^{2}=[\frac{1}{x-1}]^{2}$

\item $f(f(x))=f(\frac{1}{x-1})=\frac{1}{\frac{1}{x-1}-1}=\frac{1}{\frac{1}{x-1}-\frac{1}{1}}=\frac{1}{\frac{1-(x-1)}{x-1}}=\frac{1}{\frac{1-x+1}{x-1}}=\frac{1}{\frac{2-x}{x-1}}=\frac{\frac{1}{1}}{\frac{2-x}{x-1}}=\frac{x-1}{2-x}$
\end{itemize}

$Dom(f)=\mathbb{R}-\left\{1\right\}$

$Dom(f)=\mathbb{R}- \left\{-1,1\right\}$

$Dom(f\circ f)=\mathbb{R}-\left\{2\right\}$
\end{frame}

\end{document}
